\documentclass[conference]{IEEEtran}
\usepackage{algorithm}
\usepackage{algorithmic}
\usepackage{amsmath}
\usepackage{booktabs}
\usepackage{graphicx}

\title{Three-Tier AI Verification for Blockchain-Native Inference:\\Execution, Optimistic Attestation, and STARK Proofs}

\author{
\IEEEauthorblockN{TJ Dunham}
\IEEEauthorblockA{ARC Network}
}

\begin{document}
\maketitle

\begin{abstract}
We present a three-tier architecture for verifying AI model inference on a Layer 1 blockchain. Tier 1 provides deterministic on-chain execution via a native precompile for small models ($\leq$20B parameters). Tier 2 uses optimistic attestation with economic bonds and fraud proofs for medium-to-large models (20B--100B+). Tier 3 employs recursive STARK proofs (Stwo Circle STARKs over the Mersenne-31 field) for cryptographic verification of inference correctness. Each tier trades off computational cost, latency, and security guarantees, allowing the protocol to match verification strength to model size and application requirements.

We implement all three tiers in ARC Chain, an 80,683-line Rust blockchain with 24 native transaction types and 1,128 tests. Tier 1 executes real neural network forward passes through a dedicated precompile with deterministic left-to-right accumulation. Tier 2 manages attestation bonds, challenge periods, and escrow via on-chain state. Tier 3 generates and verifies real STARK proofs using a 32-column AIR with 9 polynomial constraints, including recursive proof composition for multi-block aggregation. Measured performance on commodity hardware (Apple M4) shows 33,230 TPS sustained with 2-round finality ($\sim$400ms), enabling high-throughput inference attestation.
\end{abstract}

\section{Introduction}

As AI models are increasingly deployed in high-stakes applications---financial trading, medical diagnosis, autonomous systems---the need to verify inference correctness becomes critical. If an AI model produces an output that affects on-chain state (e.g., liquidating a position, approving a transaction, coordinating autonomous agents), participants must be able to trust that the output was computed correctly from the claimed model and input.

No single verification approach is optimal for all scenarios. Full re-execution provides the strongest guarantee but is infeasible for large models. Zero-knowledge proofs provide cryptographic certainty but are computationally expensive---the largest model proven in a ZK circuit to date is approximately 10M parameters \cite{ezkl}, 5,000$\times$ smaller than current frontier models. Optimistic verification is efficient but introduces latency and assumes honest watchers.

We propose a three-tier architecture that combines all three approaches, automatically routing inference verification to the appropriate tier based on model size and application requirements:

\begin{enumerate}
    \item \textbf{Tier 1 --- Deterministic Execution}: For models $\leq$20B parameters, every validator executes the inference as part of normal block processing via a native precompile (address \texttt{0x0A}). The precompile calls a real neural network runtime with deterministic forward passes.

    \item \textbf{Tier 2 --- Optimistic Attestation}: For larger models (20B--100B+), a provider executes inference off-chain and submits an \texttt{InferenceAttestation} transaction containing the output hash and a bond. A challenge period allows any validator to dispute the result by posting a competing output hash with their own bond. Resolution occurs via on-chain re-execution or VRF committee vote.

    \item \textbf{Tier 3 --- STARK Verification}: For applications requiring cryptographic proof, the inference computation (or a batch of inferences) is proven using recursive Circle STARKs over the Mersenne-31 field. The proof is submitted via a \texttt{ShardProof} transaction and verified on-chain.
\end{enumerate}

This three-tier design mirrors the real-world spectrum of trust requirements: Tier 1 for routine operations where re-execution is cheap, Tier 2 for large-scale inference where economic security suffices, and Tier 3 for regulated or adversarial environments requiring mathematical proof.

\section{Architecture}

\subsection{Tier 1: On-Chain Deterministic Execution}

Tier 1 inference is executed by every validator as part of block processing. The inference engine is accessed via precompile \texttt{0x0A}, which accepts a 32-byte model identifier and input data, and returns the inference output.

\subsubsection{Neural Network Runtime}

The runtime supports five layer types: Dense (matrix multiplication with bias), ReLU, Softmax (numerically stable), LayerNorm (with learned affine parameters), and Embedding (token lookup table). Models are serialized in a binary format with little-endian \texttt{f32} weights.

\textbf{Determinism.} All matrix multiplications use strictly left-to-right sequential accumulation:

\begin{algorithmic}[1]
\FOR{$i = 0$ to $\text{out\_size} - 1$}
    \STATE $\text{acc} \leftarrow \text{bias}[i]$
    \FOR{$j = 0$ to $\text{in\_size} - 1$}
        \STATE $\text{acc} \leftarrow \text{acc} + W[i][j] \times x[j]$
    \ENDFOR
    \STATE $y[i] \leftarrow \text{acc}$
\ENDFOR
\end{algorithmic}

This avoids non-associative reductions (e.g., parallel summation with different groupings) that can produce different results on different hardware due to IEEE-754 rounding.

\subsubsection{Content-Addressed Models}

Models are identified by \texttt{model\_id = BLAKE3(weights)}. This ensures every validator loads identical weights. The identifier is committed on-chain when a model is registered, creating a tamper-evident record.

\subsection{Tier 2: Optimistic Attestation with Fraud Proofs}

For models too large for every validator to execute, Tier 2 uses an optimistic approach:

\begin{enumerate}
    \item An inference provider runs the model off-chain (on any hardware).
    \item The provider submits an \texttt{InferenceAttestation} transaction:
    \begin{itemize}
        \item \texttt{model\_id}: BLAKE3 hash of model weights
        \item \texttt{input\_hash}: BLAKE3 hash of input data
        \item \texttt{output\_hash}: BLAKE3 hash of output data
        \item \texttt{bond}: Staked collateral (returned if unchallenged)
        \item \texttt{challenge\_period}: Number of blocks for disputes
    \end{itemize}
    \item The bond is locked in a deterministic escrow account: \texttt{addr = BLAKE3("arc-inference" $\|$ tx\_hash)}.
    \item If no \texttt{InferenceChallenge} is submitted within the challenge period, the attestation is finalized and the bond is returned.
    \item If challenged, a challenger posts their own \texttt{output\_hash} with a bond. Resolution determines the winner; the loser's bond is slashed.
\end{enumerate}

\subsubsection{Economic Security}

The minimum bond must exceed the potential gain from a fraudulent attestation. If the value at stake from the inference result is $V$ and the bond is $B$, an attacker profits only if they can corrupt the result \emph{and} no honest party challenges within the period. With an active validator set monitoring attestations, the expected cost of fraud is $B$ (certain loss if caught) versus the uncertain gain of $V$ (only if uncaught). Setting $B > V$ makes fraud economically irrational regardless of the detection probability.

\subsection{Tier 3: STARK-Verified Inference}

For applications requiring cryptographic proof---regulatory compliance, high-value financial decisions, adversarial environments---Tier 3 generates a STARK proof of inference correctness.

\subsubsection{Circle STARKs over Mersenne-31}

Our STARK prover uses the Stwo framework, which implements Circle STARKs over the Mersenne-31 field ($p = 2^{31} - 1$). This provides several advantages:
\begin{itemize}
    \item The M31 field supports efficient NTT (Number Theoretic Transform) for polynomial evaluation.
    \item Circle STARKs avoid the need for a multiplicative subgroup, using circle groups instead.
    \item SIMD-friendly: 32-bit field elements pack well into 128/256/512-bit SIMD registers.
\end{itemize}

\subsubsection{Block Witness AIR}

We define an Algebraic Intermediate Representation (AIR) for block state transitions with 32 trace columns:
\begin{itemize}
    \item 1 active flag column
    \item 16 hash limb columns (BLAKE3 state hash decomposed into M31 elements)
    \item 15 transfer data columns (sender, receiver, amount, nonce, etc.)
\end{itemize}

The AIR enforces 9 constraint groups:
\begin{enumerate}
    \item Boolean constraints: active flag and transfer flag are $\{0, 1\}$.
    \item Padding constraints: inactive rows have zero data columns.
    \item Balance conservation: debits equal credits for each transfer.
    \item Nonce monotonicity: sender nonce increments by 1.
    \item Hash limb decomposition: state hash is correctly decomposed over M31.
\end{enumerate}

All constraints have degree $\leq 2$, enabling efficient proving.

\subsubsection{Recursive Composition}

For multi-block aggregation, we define a recursive verifier AIR with 82 columns:
\begin{itemize}
    \item 16 child hash columns
    \item 16 start state columns + 16 end state columns
    \item 16 Merkle sibling columns + 16 computed Merkle columns
    \item 1 chain validity column + 1 active flag
\end{itemize}

This verifier proves that a set of child block proofs are valid and that the state transitions chain correctly. The recursive structure allows logarithmic aggregation: $n$ block proofs can be aggregated into a single proof in $O(\log n)$ recursive steps.

\subsubsection{Proof Generation and Verification}

\begin{algorithmic}[1]
\STATE Generate execution trace from block state diffs
\STATE Compute polynomial commitment (Merkle tree over FFT evaluation)
\STATE Generate STARK proof via Stwo prover (SimdBackend)
\STATE Inline verification to ensure proof correctness
\STATE Serialize proof receipt: commitment roots + binding hash
\end{algorithmic}

The binding hash is $\text{BLAKE3}(\text{commitments} \| \text{block\_data})$, which ties the proof to the specific block it proves. Verification checks that this binding hash is consistent.

\section{Integration}

The three tiers integrate into a single transaction processing pipeline:

\begin{enumerate}
    \item When a transaction triggers inference (via precompile \texttt{0x0A}), the model's tier determines the verification path.
    \item Tier 1: The inference engine executes the model in the transaction's execution context. The output is available to subsequent operations in the same transaction.
    \item Tier 2: The \texttt{InferenceAttestation} transaction is processed in one block. The challenge period spans multiple blocks. Finalization occurs automatically when the period expires.
    \item Tier 3: A \texttt{ShardProof} transaction contains the STARK proof bytes. The proof is verified against the registered circuit and the claimed state transition.
\end{enumerate}

Applications can choose their tier based on requirements:
\begin{itemize}
    \item A chatbot response: Tier 2 (fast, economically secured).
    \item A financial decision affecting \$1M+: Tier 3 (cryptographic proof).
    \item A classification/scoring task with a small model: Tier 1 (instant, deterministic).
\end{itemize}

\section{Evaluation}

\subsection{Blockchain Performance}

\begin{table}[h]
\caption{Measured network performance}
\label{tab:perf}
\centering
\begin{tabular}{ll}
\toprule
Metric & Value \\
\midrule
Sustained TPS (2-node, real QUIC + DAG) & 33,230 \\
Peak TPS (single-node, M2 Ultra 24c) & 183,000 \\
Finality (2-round DAG commit) & $\sim$400 ms \\
Transaction types & 24 native \\
Codebase & 80,683 LOC Rust \\
Tests & 1,128 passing \\
\bottomrule
\end{tabular}
\end{table}

\subsection{Tier 1: Precompile Execution}

The inference precompile executes the neural network forward pass within the transaction's gas limit. For a 4-layer MLP (Dense $\rightarrow$ ReLU $\rightarrow$ Dense $\rightarrow$ Softmax):

\begin{table}[h]
\caption{Tier 1 precompile latency (per inference call)}
\centering
\begin{tabular}{ll}
\toprule
Model Size & Latency \\
\midrule
1K parameters & $< 1$ ms \\
100K parameters & $\sim$5 ms \\
1M parameters & $\sim$50 ms \\
\bottomrule
\end{tabular}
\end{table}

\subsection{Tier 2: Attestation Throughput}

Each \texttt{InferenceAttestation} is one standard transaction. At 33,230 TPS, the network can process over 33,000 attestations per second. The chain does not care about the model size because only the output hash goes on-chain.

\subsection{Tier 3: STARK Proof Performance}

Proof generation uses the Stwo SimdBackend with BLAKE3 Merkle commitments:

\begin{table}[h]
\caption{STARK proof generation (Apple M4, 10 cores)}
\centering
\begin{tabular}{lll}
\toprule
State Diffs & Proof Size & Proving Time \\
\midrule
8 & $\sim$2 KB & $\sim$15 ms \\
64 & $\sim$4 KB & $\sim$120 ms \\
256 & $\sim$8 KB & $\sim$450 ms \\
\bottomrule
\end{tabular}
\end{table}

Recursive aggregation composes multiple block proofs into a single proof. Two child proofs aggregate in $\sim$200ms on the same hardware.

\section{Related Work}

\textbf{On-chain inference.} The ``Proof of Quality'' protocol \cite{pqml} demonstrates real NLP inference on blockchains for models like Llama 3, using sampling-based verification. Our Tier 1 provides full deterministic re-execution, while our Tier 2 offers optimistic verification for larger models.

\textbf{ZK proofs of ML.} EZKL \cite{ezkl} converts ML models to ZK circuits, currently supporting models up to $\sim$10M parameters. Modulus Labs \cite{modulus} targets similar scale. Our Tier 3 uses STARK proofs rather than SNARKs, providing transparency (no trusted setup) and post-quantum security.

\textbf{Optimistic rollups for AI.} ORA \cite{ora} and Ritual \cite{ritual} use optimistic verification for off-chain AI computation. Our Tier 2 follows a similar design but is native to the L1 consensus, not a rollup, enabling tighter integration with state transitions.

\textbf{Multi-tier verification.} To our knowledge, no prior system combines deterministic execution, optimistic attestation, and STARK proofs in a single integrated architecture. Existing approaches treat these as alternatives; we treat them as complementary tiers.

\section{Conclusion}

We presented a three-tier AI verification architecture that provides deterministic execution for small models, economic security for medium-to-large models, and cryptographic proof for high-assurance applications. The architecture is implemented and tested in a production-grade L1 blockchain, demonstrating that multi-tier AI verification is both practical and performant.

The key contribution is the integration: rather than choosing a single verification strategy, applications select the appropriate tier based on model size and trust requirements. This mirrors how traditional systems combine verification mechanisms (checksums for routine data, digital signatures for important data, zero-knowledge proofs for privacy-sensitive data) rather than applying the strongest guarantee universally.

\begin{thebibliography}{9}
\bibitem{ezkl} J. Cheng et al., ``EZKL: Zero-knowledge proofs for machine learning,'' 2023.
\bibitem{modulus} Modulus Labs, ``The cost of intelligence: Proving ML inference on-chain,'' 2023.
\bibitem{pqml} ``Proof of Quality: A costless paradigm for trustless generative AI model inference on blockchains,'' arXiv:2405.17934, 2024.
\bibitem{ora} ORA Protocol, ``Optimistic machine learning on Ethereum,'' 2024.
\bibitem{ritual} Ritual, ``Infernet: Decentralized AI inference network,'' 2024.
\bibitem{stwo} StarkWare, ``Stwo: Circle STARKs over M31,'' 2024.
\end{thebibliography}

\end{document}
